\chapter{Data and Descriptive Evidence}
\label{chap:data}

\section{Data sources}

The analysis combines two datasets from the Varieties of Democracy project. The V-Dem Country-Year
Dataset v16 provides annual political indicators, including corruption, electoral democracy, income,
population, and region \citep{vdem2026}. The V-Party V2 dataset provides party-level expert assessments
of ideology and organisation at election dates, including populism and government status
\citep{vparty2022}. V-Dem's measurement model aggregates multiple expert ratings into comparable
latent estimates while explicitly modelling coder reliability and uncertainty
\citep{pemstein2026}. These data offer unusually broad temporal and geographic coverage, but they remain
expert-coded measures rather than direct administrative records.

The two sources have different units of observation. V-Dem is a country-year panel. V-Party is a
country-party-election dataset and is empirically dense only in election years. Linking them therefore
requires substantive choices about which party represents the government and how an election-year
score should be carried into subsequent years. \Cref{fig:design} summarises the complete workflow.

\begin{figure}[H]
  \centering
  \includegraphics[width=\textwidth]{research_design.png}
  \caption{Research design and analytical workflow. Party-election populism scores are filtered,
  aggregated, forward-filled, and joined to annual V-Dem outcomes before the three research questions
  are estimated.}
  \label{fig:design}
\end{figure}

\section{Core measures}

\subsection{Governing-party populism}

The primary explanatory variable is V-Party's \codevar{v2xpa\_popul}, a continuous index from 0 to 1
constructed from people-centrism and anti-elitism. This operationalisation matches the ideational
definition of populism as a thin-centred opposition between the people and the elite
\citep{mudde2004}. The index measures rhetoric and party identity; it should not be read as a complete
measure of nationalism, policy radicalism, or democratic quality.

The analysis focuses on parties coded as the senior governing partner
(\codevar{v2pagovsup = 0}), normally the prime minister's party. If a country-election has more than one
senior governing party, their populism scores are averaged using vote shares when all weights are
available and a simple mean otherwise. Populism is modelled continuously. The value 0.47, the 75th
percentile in the raw governing-party distribution, is used only for descriptive labels and never as a
regression threshold.

\subsection{Political corruption}

The primary outcome, \codevar{v2x\_corr}, is V-Dem's composite political corruption index. It ranges
from 0 (cleaner) to 1 (more corrupt) and averages corruption in four institutional spheres. The
decomposition uses executive corruption (\codevar{v2x\_execorr}), public-sector corruption
(\codevar{v2x\_pubcorr}), legislative corruption (\codevar{v2lgcrrpt}), and judicial corruption
(\codevar{v2jucorrdc}). The last two continuous latent estimates use the opposite source orientation:
higher values indicate cleaner institutions. They are multiplied by $-1$ in decomposition models so
that a positive coefficient always means more corruption.

The project also constructs five-level 0--1 versions of legislative and judicial corruption from the
ordinal ratings. These harmonised variables are useful for plots but are too coarse to be the primary
fixed-effects outcomes: several countries never change ordinal category. The continuous latent
estimates preserve substantially more within-country information.

\begin{table}[H]
\centering
\caption{Variables used in the empirical analysis.}
\label{tab:variables}
\begin{threeparttable}
\small
\begin{tabularx}{\textwidth}{>{\raggedright\arraybackslash}p{3.1cm}>{\raggedright\arraybackslash}p{3.0cm}Y>{\raggedright\arraybackslash}p{2.5cm}}
\toprule
Role & Variable & Interpretation & Model use \\
\midrule
Treatment & \texttt{populism\_}\newline\texttt{governing} & Senior governing party's continuous populism score,
forward-filled between elections & All primary models \\
Composite outcome & \codevar{v2x\_corr} & Political corruption, 0 cleaner to 1 more corrupt & RQ1--RQ2 \\
Executive outcome & \codevar{v2x\_execorr} & Bribery and embezzlement by the executive & RQ3 \\
Public-sector outcome & \codevar{v2x\_pubcorr} & Bribery and theft in public administration & RQ3 \\
Legislative outcome & \codevar{v2lgcrrpt} & Continuous legislative corruption estimate; sign reversed in models & RQ3 \\
Judicial outcome & \codevar{v2jucorrdc} & Continuous judicial corruption estimate; sign reversed in models & RQ3 \\
Economic control & \codevar{log\_gdppc} & Natural logarithm of GDP per capita & FE and direction models \\
Political control & \codevar{v2x\_polyarchy} & Electoral democracy, 0 to 1 & FE and direction models \\
Context & \codevar{e\_regiongeo}, \codevar{e\_pop} & Region and population & Description / CV only \\
\bottomrule
\end{tabularx}
\begin{tablenotes}[flushleft]\footnotesize
\item Region is time-invariant and therefore absorbed by country fixed effects. Population is not an
estimated control. This avoids presenting columns retained for description as if they entered the
reported models.
\end{tablenotes}
\end{threeparttable}
\end{table}

\section{Panel construction}

The analytical window is 1970--2019. V-Party contains records before 1970, but its expert-coded
populism and government-status variables are unpopulated in those earlier years. Coverage ends in
2019, so the analysis does not include the most recent populist-government episodes.

The country-year panel is built in five steps:

\begin{enumerate}
  \item Restrict V-Party to senior governing parties and the 1970--2019 window.
  \item Require more than three coders for both expert-coded components of the populism index. The
  composite index itself has no coder-count field, so filtering its components is the appropriate
  quality check.
  \item Aggregate multiple senior governing parties to one country-election observation.
  \item Forward-fill the election observation until the next coded election, without filling years
  before a country's first observation.
  \item Inner-join the result to V-Dem on \codevar{country\_id} and \codevar{year}; validate unique,
  gap-free annual histories before constructing lags and leads.
\end{enumerate}

Forward-filling is necessary to align election-dated party information with annual corruption data,
but it is also an assumption: governing-party populism is treated as constant between elections. The
models therefore contain many annual rows but fewer genuinely distinct changes in government
populism. Election-year and government-change-year estimates are reported as robustness checks.

\begin{table}[H]
\centering
\caption{Sample construction and effective estimation samples.}
\label{tab:sample}
\small
\begin{tabularx}{0.9\textwidth}{Yrrl}
\toprule
Stage & Observations & Countries & Interpretation \\
\midrule
Senior governing-party observations before strict coder filter & 1,628 & 163 & Election-dated source \\
After component coder-count filter & 953 & 96 & Reliable election observations \\
Final annual panel after forward-fill and join & 3,965 & 96 & 1970--2019 panel \\
M1--M3 complete-case stability sample & 3,940 & 95 & Main association sample \\
M4 sample with lagged corruption & 3,845 & 95 & Dynamic short-run model \\
Two-lag composite direction sample & 3,750 & 95 & Main lead-lag / Granger sample \\
Legislative outcome available & 3,770 & 95 & Structurally lower coverage \\
\bottomrule
\end{tabularx}
\end{table}

The strict coder filter trades representation for reliability. It removes 67 countries present before
filtering, and the resulting estimates should not be generalised automatically to countries with
sparser expert coverage. At the same time, the panel retains 95--96 countries and up to 50 annual
observations per country, which is substantial for a within-country design.

\section{Exploratory evidence}

\subsection{Distribution and trends}

In the broader election-year source, governing-party populism is unimodal and right-skewed
(mean 0.330, median 0.274, 75th percentile 0.470). There is no empirical break that would justify
turning the index into a populist/non-populist dummy. The global mean drifts upward from roughly 0.30
in the 1970s to 0.36 in the 2010s, but the pattern is gradual and noisy rather than a single post-2010
break.

Composite corruption is bimodal, with one cluster of clean democracies near zero and another of
high-corruption regimes near 0.85. Regional differences are much larger than the descriptive
high-versus-lower-populism gap: mean corruption ranges from approximately 0.05 in Western Europe and
North America to approximately 0.85 in Central Asia. These differences motivate the separation of
between-country and within-country variation.

\begin{figure}[H]
  \centering
  \includegraphics[width=0.95\textwidth]{eda_distributions.png}
  \caption{Distributions of the treatment and composite outcome in the strict modelling panel. Dashed
  lines mark medians. These histograms describe the estimation sample and are not evidence of a causal
  relationship.}
  \label{fig:distributions}
\end{figure}

\subsection{Descriptive group contrasts}

For presentation only, the exploratory notebook labels the top quartile of governing-party populism
as ``high populism.'' Across 1,831 high-populism and 5,078 lower-populism country-years in the broader
descriptive panel, all corruption measures are higher in the high-populism group. Yet Cliff's delta is
small for every dimension (approximately 0.05--0.11). Legislative and judicial corruption show the
largest gaps, whereas executive and public-sector corruption show the smallest. Large sample size
makes the group differences statistically detectable; it does not make them substantively large.

This descriptive result is vulnerable to confounding because it compares different countries and
regions. It is therefore treated as motivation for RQ1 and RQ3, not as a finding about what happens
after a country elects a more populist government.

\subsection{Illustrative governing transitions}

Five case studies reinforce the ambiguity. Comparing five-year average composite corruption before
and after a documented populist entry, Venezuela increases from 0.69 to 0.87 and Hungary from 0.28 to
0.40. Turkey changes only slightly (0.50 to 0.53), Bolivia is essentially flat (0.70 to 0.70), and
Ecuador declines slightly (0.71 to 0.67). Venezuela, Bolivia, and Ecuador were already relatively
corrupt before populist entry, whereas Hungary begins from a cleaner baseline and more closely
resembles an institutional-capture narrative.

The cases demonstrate why neither a before/after anecdote nor a cross-sectional correlation can
resolve direction. Populism may precede corruption, corruption may create demand for populism, or both
may reflect common institutional and economic shocks. The next chapter translates these alternatives
into distinct panel estimands.

\section{Measurement and data-quality risks}

Four limitations are carried into the empirical design. First, V-Dem and V-Party are both
expert-coded, so shared perceptions can induce correlated measurement error. Second, forward-filled
populism changes mainly around elections, while corruption is extremely persistent. Third,
legislative corruption is missing when no legislature exists, making missingness institutionally
meaningful rather than random. Fourth, the five-level legislative and judicial ordinal measures are
often constant within countries. The report therefore uses continuous point estimates for primary
models, clustered inference, first-difference checks, election-based samples, and explicit diagnostics
of within-country variation.
